\subsection{The manuscript's statements and relative class numbers}

\begin{claim}{C.R1}{Proved}{The two curve assertions in cote 149, scan 36}
Over every field finitely generated over its prime field, the
point-section map for a smooth hyperbolic curve is injective:
use the full geometric group in characteristic zero, and its
prime-to-$p$ quotient in characteristic $p$. For a proper curve
of genus at least two over a finite field, the covering condition
of scan 36 forces $P=Q$.
\end{claim}
\begin{proof}
In positive characteristic this is exactly
\statementref{A.R1} and \statementref{A.R5}. Their dependencies
are the same-prime compactification lemma, the cuspidal criterion
with its orientation and Kummer evaluation, the relative
configuration exactness, and the diagonal boundary transport;
each was independently proved in Part A. The covering condition
is equivalent to section conjugacy by \statementref{A.R3}, and
the main theorem then gives equality. Characteristic zero is the
all-prime Kummer argument in \statementref{A.R2} and
\statementref{A.R5}.

Thus the sentence in scan 36, \emph{«en déduire $P=Q$»}, holds in
the specified curve scope. The phrase
\emph{«sans hyp. de caractéristique»} requires retaining the
finite-generation and hyperbolicity hypotheses. It does not make
the corresponding abelianized assertion true in positive
characteristic, and it does not include proper elliptic curves.
\end{proof}

\begin{claim}{C.R2}{Refuted}{The characteristic-$p$ prime-to-$p$ extension of Théorème 3 in all dimensions}
The asserted characteristic-$p$ extension of cote 149, scan 32,
Théorème 3, for every geometrically connected finite-type scheme
embedded in a semiabelian variety, is false, already in dimension
one. Over $\mathbf F_2$, take
\[
 E:\quad y^2+xy=x^3+1,\qquad O=\infty,\quad T=(0,1).
\]
The two points have identical prime-to-two section classes in
the whole arithmetic extension. For an arbitrary semiabelian
$G/K$ over a finitely generated characteristic-$p$ field, the
exact replacement is
\[
 \ker\bigl(G(K)\to\hbox{prime-to-}p\hbox{ section classes}\bigr)
       =G(K)[p^\infty]
\]
after using the identity point to identify this section set with
abelian Galois cohomology. A subscheme embedded in $G$ has
injective point sections whenever no difference of distinct
rational points lies in $G(K)[p^\infty]$.
\end{claim}
\begin{proof}
The explicit checks in \statementref{A.CR11} verify that $E$ is
smooth and $T$ has exact order two. Its geometric fundamental
group is abelian, its prime-to-two part is the product of its
prime-to-two Tate modules, and the difference of the two point
sections is the Kummer class of $T$, which is zero. This is a
conjugacy assertion in the whole group, not a finite-level test.

For a semiabelian variety the same abelian fundamental-group
description and \statementref{A.R2} give the exact kernel.
Functoriality of fundamental groups for an embedding $X\to G$
shows that a collision on $X$ maps to a point difference in this
kernel. Its avoidance is therefore a sufficient condition on
$X$. We do not assert the converse for $X$: hyperbolic curves
may still separate points nonabelianly when their generalized
Jacobian Kummer map fails to do so, as the negative control
\statementref{A.CR12} and the main theorem show. Products
$E\times\mathbf G_m^{d-1}$ give the same counterexample in every
positive dimension.

This matches the warning on scan 34 about translation by
$p$-torsion and the failure of the abelianized map (**) on scan
35. Removing a point from an elliptic curve can restore the
nonabelian hyperbolic assertion, as proved in Part A; it does
not invalidate the warning about (**).
\end{proof}

\paragraph{Readings and the limits of the transcript check.}
We checked the supplied second-pass transcriptions, not absent
images of the original scans. The alternatives
\emph{à l'aise / l'air} on scan 35 do not affect a hypothesis.
The margin's uncertain reading involving an \emph{anneau} is
mathematically consistent with spreading to a finite-type model
and specializing at a closed point, as in \statementref{A.R4};
the proof does not depend on selecting that reading. The
alternative \emph{ou / en} in the preceding sentence also changes
no mathematical quantifier used here. In scan 36, a cover $Y'$
in the displayed covering condition is a cover of $Y$, as its
chosen point and generalized Jacobian require; the transcription's
$X$ is the noted variable slip. We interpret the torsion there
as $p$-primary, rather than as necessarily killed by $p$.
No statement in this paper depends on resolving a doubtful word.

\begin{claim}{C.R3}{Proved}{The precise class-number consequences and a sharp limitation}
Point separation gives the isolating pointed cover of
\statementref{B.R1}; it gives no numerical value for the
$p$-part of its relative class number. There is no universal
assertion that every nontrivial pointed prime-to-$p$ cover has
$p$-unit relative class number, or that every such cover has
$p$-divisible relative class number. Both behaviours occur for
the same curve, point, field, and prime degree in our computation.

For $C_\alpha/\mathbf F_{64}$ of \statementref{B.R3},
$\alpha$ with bit code $58$, and $P=\infty$, the degree-seven
characters $(1,0)$ and $(1,1)$ in the recorded basis give
respectively
\[
 h(D_{(1,0)})/h(C)=46713473689,\qquad
 h(D_{(1,1)})/h(C)=55586549824.
\]
The first is odd and the second is even. Nevertheless the
point-section map of $C$ is injective over this field and every
constant extension.
\end{claim}
\begin{proof}
The exact occupancy vectors for these two characters are
$(9,12,17,12,12,17,12)$ and $(11,18,14,8,8,14,18)$.
The degree-two twisted polynomials, including leading coefficient
64 and the infinity pointing, are supplied in
\statementref{B.R3}. Their full integer resultants give exactly
the two displayed integers; both are retained in the computed
JSON, and neither is a conjectural value. The main theorem and
its checked inputs apply to this same curve.

A hypothetical collision would force the relative divisibility
in the applicable class-number criterion. Its nonexistence
removes that premise; contraposition does not turn it into
existence of a cover with nondivisible relative class number.
Our explicit constructions prove such existence for the
fourteen requested $\mathbf F_4$ pointed cases and for the
supersingular test curve at primes 191 and 229, as shown in
Part B. Those are additional calculations, not formal
consequences of the point-separation theorem. The universal
quantifier in B1 therefore remains a separate question.
\end{proof}

\begin{claim}{C.R4}{Proved}{What transfers to the other research line}
The characteristic-zero construction of a virtual puncture
does not prove the surjectivity of the section map for the
tripod over $\mathbf Q$, and does not construct a birational
lift of an arbitrary section. The present finite-field
injectivity theorem supplies no such assertion. Likewise a
non-point-theoretic full section over a finite field obtained
from a norm-module class in Part B is not a counterexample
over a field finitely generated over $\mathbf Q$.
\end{claim}
\begin{proof}
The pullback construction in the next subsection supplies an
extension for every section, without giving a rational point or
a splitting of its virtual cusp. Recognizing it as the punctured
group of an actual point is an additional condition. The
characteristic-zero separation criterion compares two already
geometric point sections; it has no arbitrary section as a
second input. The finite-field norm-module realization changes
both the characteristic and the arithmetic Galois group, so no
base-change map transfers it to $\mathbf Q$. These are exact
limitations of the maps constructed, not assertions that the
characteristic-zero surjectivity or birational-lifting questions
are false. They address the separate LM6 line referred to in
the supplied intake.
\end{proof}

